\section{System and Simulation Model}
\label{sec:model}

A workflow is a directed acyclic graph (DAG) of model calls and tool delays.
Completion time runs from its first call's release until all required calls
finish. Loss on a parallel branch can be absorbed by work elsewhere, whereas
loss on a critical branch can delay completion. Protection value therefore
depends on residual workflow completion.

\paragraph{Execution scope.}
The CPU simulator uses one global virtual clock for admission, scheduler
steps, reservations, transfers, failure and repair. Eight DP workers each
hold Qwen3-8B with TP=1 and PP=1. Per-GPU HBM ledgers account for committed
cache, active primary KV, pending reservations, staging, weights and runtime
reserve. Store-DRAM has separate capacity and transfer accounting. One
colocated Store uses PCIe paths; SSD/NVMe and remote Store are excluded.

\paragraph{Continuous batching and timing.}
The scheduler-step bridge models residency, preemption and transport
contention without executing CUDA kernels. It distributes each request's
ideal duration over token work, with prefill weight 1 and decode weight 8.
A shared step takes the maximum participant share. This ideal parallel
model omits shape-dependent kernel efficiency and hardware interference.
Admission order changes contention and routing, while graphs and token
content remain fixed.
For assumption sensitivity, we also use
$d_{\beta}=\lceil(1-\beta)\max_i d_i+\beta\sum_i d_i\rceil$ with
$\beta\in\{0,0.5,1\}$. Execution and routing's queue estimate use the same
composition. These are alternative service assumptions, not physical latency
bounds; the single-request service budget remains the same.

\paragraph{Failure and repair.}
A GPU failure fences the device and destroys its HBM. Recovery uses
committed surviving prefixes, recomputes missing prompt tokens and silently
replays prior output before resuming visible decoding. Generation and route
checks fence stale work. Repair restores empty capacity; commit or abort
releases reservations and source pins. Pending copies cannot recover state.

\section{AegisKV Policy}
\label{sec:policy}

The source worker independently admits sealed chunks to Peer-HBM and
Store-DRAM using residual slack and a feedback-controlled Peer envelope.
Peer protection does not depend on a completed Store copy, and eviction
does not destage between tiers. Reservations, transfers and commitments
are independent.

\paragraph{Algorithm 1: ideal residual slack.}
At causal request and tool boundaries, the simulator computes an exact critical
path on the residual DAG using the workload's ideal call durations. For call
$i$, let $s_i(t_0)$ be published slack at time $t_0$. Between observations,
\begin{equation}
  s_i(t)=\max\{0,s_i(t_0)-(t-t_0)\}.
\end{equation}
This ideal-duration experiment uses no GBDT predictor or prediction error.
Slack uses neither future failures nor observed absolute finish times;
the findings therefore do not validate online prediction.

\paragraph{Algorithm 2: tier-local protection value.}
For request $i$ and pool $p$, let $T_{i,p}$ be modeled recovery time given
eligible prefixes, including loading, missing prompt computation, and silent
replay. Define
\begin{equation}
  J_{i,p}(t)=\max\{T_{i,p}(t)-s_i(t),0\}.
\end{equation}
Candidate gain is
$A_c=[J_{i,p}^{\mathrm{before}}-J_{i,p}^{\mathrm{after}}]^+$.
The evaluated policy normalizes recovery-harm reduction by the physical
bytes allocated for the candidate.

Specifically, pool $p$ maintains rent $L_p$. At commitment, chunk $c$ with
current gain $A_c$ and physical size $b_c$ receives key
$H_c=L_p+A_c/b_c$ and retains residual credit
$b_c[H_c-L_p]^+$. Only the currently safe suffix tail of each request enters
the pool's minimum-key heap. When its deletion loss changes from
$D_c^{\mathrm{old}}$ to $D_c^{\mathrm{new}}$, its key becomes
\begin{equation}
 H_c\leftarrow\max\!\left\{L_p,
 H_c+\frac{D_c^{\mathrm{new}}-D_c^{\mathrm{old}}}{b_c}\right\}.
\end{equation}
Hidden chunks preserve their previous keys and refresh when exposed as tails;
exposure does not grant fresh credit at the current rent.

When space is needed, the landlord previews a suffix-safe victim set $V$.
Loss $D(V)$ is recomputed jointly, once per affected request, from tentative
prefix state. Summing independent chunk credits could count lost prefix
utility repeatedly. Admission requires every selected key to be below
$\Theta_c=L_p+A_c/b_c$ and $A_c>D(V)$; a tie is rejected.
After eviction, rent advances to the larger of its current value and the
largest selected key. The new chunk reserves capacity in \texttt{PENDING}
state and refreshes gain and rent when committing its key. A failed transfer
releases its reservation without restoring the evicted copies. During this
replacement window, recovery uses other committed prefixes or recomputation.
Preview excludes dependencies needed to realize
the candidate's gain and accounts for physical aliases. Peer-HBM and
Store-DRAM retain separate rents, capacities, reservations, and victims.

Candidate fractions describe overlapping prefixes starting at token zero.
Default Peer and Store candidates cover 75\% and 90\%, respectively; these
percentages are not additive coverage. Rejection, capacity, incomplete transfer,
eviction, and failure can reduce committed coverage below this envelope.

\paragraph{Routing and recovery-source choice.}
The cache-aware router selects a serving GPU by estimated completion time,
combining an eligible continuous prefix, current queue state, and estimated
restore time. After a failure destination is chosen, the recovery selector
compares Peer-first and Store-first ownership of overlapping copies. Both
assignments use normal HBM first and preserve the same effective prefix and
recompute suffix. Among feasible assignments, the lexicographic key is
estimated completion, admission, materialization, transferred bytes, and
Peer-first as the exact-tie winner. Candidate evaluation is side-effect free, using active flows, proposed
transfers and current destination load, without predicting future arrivals
or batches. Both estimates are retained, but only the selected assignment
executes. No permanently Peer-first arm isolates this selector's benefit.
Source substitution preserves recomputation when the effective prefix is
unchanged.

\paragraph{Algorithm 3: Peer-HBM quota feedback.}
The initial envelope is $\alpha M$, where $M$ is KV-capable HBM after weights
and fixed runtime reserve. At fixed control-window boundaries, foreground
pressure causes $\alpha\leftarrow\max(\alpha_{\min},\gamma\alpha)$.
After $K$ stable windows,
$\alpha\leftarrow\min(\alpha_{\max},\alpha+\delta)$.
Quota is $\lfloor\alpha M\rfloor$ bytes. Reduction can block admission while
safe tails are reclaimed; it cannot free pinned or active state by changing
an accounting number. Foreground allocation and recovery still require
physical capacity transactions. The default window is the larger of one
second and the median ideal workflow CPM makespan.
An unavailable GPU's control window clears old pressure without changing
quota; post-repair decisions use pressure observed in subsequent windows.

\paragraph{Executable checks.}
Runs retain scheduler steps, residency, physical ledgers, transfer events,
policy previews, and quota transitions. Offline checks reconstruct completion,
exact rational gain and loss, byte and rent conservation, pending-copy
invisibility, and AIMD equations. These establish simulation consistency,
not physical accuracy of the latency model.
